\documentclass[12pt,letterpaper]{article}
\usepackage{graphicx,textcomp}
\usepackage{natbib}
\usepackage{setspace}
\usepackage[margin=1in]{geometry}
\usepackage{color}
\usepackage[reqno]{amsmath}
\usepackage{amsthm}
\usepackage{fancyvrb}
\usepackage{amssymb,enumerate}
\usepackage[all]{xy}
\usepackage{endnotes}
\usepackage{lscape}
\newtheorem{com}{Comment}
\usepackage{float}
\usepackage{hyperref}
\newtheorem{lem} {Lemma}
\newtheorem{prop}{Proposition}
\newtheorem{thm}{Theorem}
\newtheorem{defn}{Definition}
\newtheorem{cor}{Corollary}
\newtheorem{obs}{Observation}
\usepackage[compact]{titlesec}
\usepackage{dcolumn}
\usepackage{tikz}
\usetikzlibrary{arrows}
\usepackage{multirow}
\usepackage{xcolor}
\newcolumntype{.}{D{.}{.}{-1}}
\newcolumntype{d}[1]{D{.}{.}{#1}}
\definecolor{light-gray}{gray}{0.65}
\usepackage{url}
\usepackage{listings}
\usepackage{color}

\definecolor{codegreen}{rgb}{0,0.6,0}
\definecolor{codegray}{rgb}{0.5,0.5,0.5}
\definecolor{codepurple}{rgb}{0.58,0,0.82}
\definecolor{backcolour}{rgb}{0.95,0.95,0.92}

\lstdefinestyle{mystyle}{
	backgroundcolor=\color{backcolour},   
	commentstyle=\color{codegreen},
	keywordstyle=\color{magenta},
	numberstyle=\tiny\color{codegray},
	stringstyle=\color{codepurple},
	basicstyle=\footnotesize,
	breakatwhitespace=false,         
	breaklines=true,                 
	captionpos=b,                    
	keepspaces=true,                 
	numbers=left,                    
	numbersep=5pt,                  
	showspaces=false,                
	showstringspaces=false,
	showtabs=false,                  
	tabsize=2
}
\lstset{style=mystyle}
\newcommand{\Sref}[1]{Section~\ref{#1}}
\newtheorem{hyp}{Hypothesis}

\title{Problem Set 3}
\date{Due: November 13, 2025}
\author{Applied Stats/Quant Methods 1}


\begin{document}
	\maketitle
	\section*{Instructions}
	\begin{itemize}
		\item Please show your work! You may lose points by simply writing in the answer. If the problem requires you to execute commands in \texttt{R}, please include the code you used to get your answers. Please also include the \texttt{.R} file that contains your code. If you are not sure if work needs to be shown for a particular problem, please ask.
	\item Your homework should be submitted electronically on GitHub.
	\item This problem set is due before 23:59 on Thursday November 13, 2025. No late assignments will be accepted.

	\end{itemize}

		\vspace{.25cm}
	
\noindent In this problem set, you will run several regressions and create an add variable plot (see the lecture slides) in \texttt{R} using the \texttt{incumbents\_subset.csv} dataset. Include all of your code.
\vspace{.5cm}
\section*{Question 1}
\vspace{.25cm}
\noindent We are interested in knowing how the difference in campaign spending between incumbent and challenger affects the incumbent's vote share. 
	\begin{enumerate}
		\item Run a regression where the outcome variable is \texttt{voteshare} and the explanatory variable is \texttt{difflog}.	
\lstinputlisting[language=R, firstline=34, lastline=41]{PS03_DP.R}
\newpage
\begin{table}[!htbp] \centering 
	\caption{} 
	\label{} 
	\begin{tabular}{@{\extracolsep{5pt}}lc} 
		\\[-1.8ex]\hline 
		\hline \\[-1.8ex] 
		& \multicolumn{1}{c}{\textit{Dependent variable:}} \\ 
		\cline{2-2} 
		\\[-1.8ex] & Incumbent Vote Share \\ 
		\hline \\[-1.8ex] 
		Campaign Spending Difference & 0.042$^{***}$ \\ 
		& (0.001) \\ 
		& \\ 
		Constant & 0.579$^{***}$ \\ 
		& (0.002) \\ 
		& \\ 
		\hline \\[-1.8ex] 
		\textit{Note:}  & \multicolumn{1}{r}{$^{*}$p$<$0.1; $^{**}$p$<$0.05; $^{***}$p$<$0.01} \\ 
	\end{tabular} 
\end{table} 
		\item Make a scatterplot of the two variables and add the regression line. 	
		\begin{figure}[h!]\centering
			\caption{\footnotesize scatterplot of difflog and voteshare}
			\label{fig:plot_1}
			\includegraphics[width=.82\textwidth]{scatter_difflog_vs.pdf}
		\end{figure} 
		\lstinputlisting[language=R, firstline=43, lastline=51]{PS03_DP.R}
		\vspace{.5cm}
		\item Save the residuals of the model in a separate object.	
		\lstinputlisting[language=R, firstline=53, lastline=55]{PS03_DP.R}
		\vspace{.5cm}
		\item Write the prediction equation. \\ 
		\begin{center} 
			$\widehat{\text {VoteShare}} = \hat{\alpha} + \hat{\beta} \times \text{DiffLog}$ \\ [0.2cm]
			$\widehat{\text {VoteShare}} = 0.579 + 0.042 \times \text{DiffLog} $
		\end{center}
		The intercept $\hat{\alpha} = 0.579$, represents the expected incumbent vote share when the difference in campaign spending between incumbent and challenger is 0 (when both spend the same amount). The slope $\hat{\beta} = 0.042$ shows that for every one-unit increase in campaign spending difference, the incumbent vote share increases by 0.042, on average. 
	\end{enumerate} 
\vspace{1.5cm}
\section*{Question 2}
\noindent We are interested in knowing how the difference between incumbent and challenger's spending and the vote share of the presidential candidate of the incumbent's party are related.	\vspace{.25cm}
	\begin{enumerate}
		\item Run a regression where the outcome variable is \texttt{presvote} and the explanatory variable is \texttt{difflog}.	
		\lstinputlisting[language=R, firstline=58, lastline=60]{PS03_DP.R}
		\newpage
\begin{table}[!htbp] \centering 
			\caption{} 
			\label{} 
		\begin{tabular}{@{\extracolsep{5pt}}lc} 
				\\[-1.8ex]\hline 
				\hline \\[-1.8ex] 
				& \multicolumn{1}{c}{\textit{Dependent variable:}} \\ 
				\cline{2-2} 
				\\[-1.8ex] & Presidential Vote Share \\ 
				\hline \\[-1.8ex] 
				Campaign Spending Difference  & 0.024$^{***}$ \\ 
				& (0.001) \\ 
				& \\ 
				Constant & 0.508$^{***}$ \\ 
				& (0.003) \\ 
				& \\ 
				\hline \\[-1.8ex] 
				\textit{Note:}  & \multicolumn{1}{r}{$^{*}$p$<$0.1; $^{**}$p$<$0.05; $^{***}$p$<$0.01} \\ 
		\end{tabular} 
\end{table} 
		\item Make a scatterplot of the two variables and add the regression line. 
		\begin{figure}[h!]\centering
			\caption{\footnotesize scatterplot of difflog and presvote}
			\label{fig:plot_2}
			\includegraphics[width=.82\textwidth]{scatter_difflog_pv.pdf}
		\end{figure} 
		\item Save the residuals of the model in a separate object.	
		\lstinputlisting[language=R, firstline=72, lastline=74]{PS03_DP.R}
		\vspace{.5cm}
		\item Write the prediction equation.
		\begin{center}
		$\widehat{\text {PresVote}} = \hat{\alpha} + \hat{\beta} \times \text{DiffLog}$ \\ [0.2cm]
		$\widehat{\text {PresVote}} = 0.508 + 0.024 \times \text{DiffLog}$
		\end{center}
		The intercept $\hat{\alpha} = 0.508$, represents the expected presidential vote share when the difference in campaign spending between incumbent and challenger is 0 (when both spend the same amount). The slope $\hat{\beta} = 0.024$ shows that for every one-unit increase in campaign spending difference, the presidential vote share increases by 0.024, on average. 
	\end{enumerate}
	
\vspace{1.5cm}
\section*{Question 3}

\noindent We are interested in knowing how the vote share of the presidential candidate of the incumbent's party is associated with the incumbent's electoral success.
	\vspace{.25cm}
	\begin{enumerate}
		\item Run a regression where the outcome variable is \texttt{voteshare} and the explanatory variable is \texttt{presvote}.
			\lstinputlisting[language=R, firstline=77, lastline=79]{PS03_DP.R}
			\begin{table}[!htbp] \centering 
				\caption{} 
				\label{} 
				\begin{tabular}{@{\extracolsep{5pt}}lc} 
					\\[-1.8ex]\hline 
					\hline \\[-1.8ex] 
					& \multicolumn{1}{c}{\textit{Dependent variable:}} \\ 
					\cline{2-2} 
					\\[-1.8ex] & Incumbent Vote Share \\ 
					\hline \\[-1.8ex] 
					Presidential Vote Share & 0.388$^{***}$ \\ 
					& (0.013) \\ 
					& \\ 
					Constant & 0.441$^{***}$ \\ 
					& (0.008) \\ 
					& \\ 
					\hline \\[-1.8ex] 
					\textit{Note:}  & \multicolumn{1}{r}{$^{*}$p$<$0.1; $^{**}$p$<$0.05; $^{***}$p$<$0.01} \\ 
				\end{tabular} 
			\end{table} 
		\newpage
		\item Make a scatterplot of the two variables and add the regression line. 
		\begin{figure}[h!]\centering
			\caption{\footnotesize Scatterplot of presvote and voteshare}
			\label{fig:plot_3}
			\includegraphics[width=.82\textwidth]{scatter_pv_vs.pdf}
		\end{figure} 
		\item Write the prediction equation.
		\begin{center}
			$\widehat{\text {VoteShare}} = \hat{\alpha} + \hat{\beta} \times \text{PresVote}$ \\ [0.2cm]
			$\widehat{\text {VoteShare}} = 0.441 + 0.388 \times \text{PresVote}$
		\end{center}
		The intercept $\hat{\alpha} = 0.441$, represents the expected incumbent vote share when the ipresidential vote share is 0. Nevertheless, since a presidential vote share of 0 is not a realistic result, the value of the intercept has little to no substantive meaning. The slope $\hat{\beta} = 0.388$ shows that for every one-unit increase in presidential vote share, the incumbent vote share increases by 0.388, on average. 
	\end{enumerate}
	

\newpage	
\section*{Question 4}
\noindent The residuals from part (a) tell us how much of the variation in \texttt{voteshare} is $not$ explained by the difference in spending between incumbent and challenger. The residuals in part (b) tell us how much of the variation in \texttt{presvote} is $not$ explained by the difference in spending between incumbent and challenger in the district.
	\begin{enumerate}
		\item Run a regression where the outcome variable is the residuals from Question 1 and the explanatory variable is the residuals from Question 2.	
		\vspace{0.5cm}
		\lstinputlisting[language=R, firstline=92, lastline=98]{PS03_DP.R}
		\begin{table}[!htbp] \centering 
			\caption{} 
			\label{} 
			\begin{tabular}{@{\extracolsep{5pt}}lc} 
				\\[-1.8ex]\hline 
				\hline \\[-1.8ex] 
				& \multicolumn{1}{c}{\textit{Dependent variable:}} \\ 
				\cline{2-2} 
				\\[-1.8ex] & Residuals 1 \\ 
				\hline \\[-1.8ex] 
				Residuals 2  & 0.257$^{***}$ \\ 
				& (0.012) \\ 
				& \\ 
				Constant & $-$0.000 \\ 
				& (0.001) \\ 
				& \\ 
				\hline \\[-1.8ex] 
				\textit{Note:}  & \multicolumn{1}{r}{$^{*}$p$<$0.1; $^{**}$p$<$0.05; $^{***}$p$<$0.01} \\ 
			\end{tabular} 
		\end{table} 
		\newpage
		\item Make a scatterplot of the two residuals and add the regression line. 
			\begin{figure}[h!]\centering
			\caption{\footnotesize scatterplot of residuals 1 and residuals 2}
			\label{fig:plot_4}
			\includegraphics[width=.82\textwidth]{scatter_res2_res1.pdf}
		\end{figure} 
		\item Write the prediction equation.
			\begin{center}
			$\widehat{\text {Residuals1}} = \hat{\alpha} + \hat{\beta} \times \text{Residuals2}$ \\ [0.2cm]
			$\widehat{\text {Residuals1}} = 0 + 0.257 \times \text{Residuals2}$ \\ [0.2cm]
			$\widehat{\text {Residuals1}} = 0.257 \times \text{Residuals2}$ 
		\end{center}
		The intercept $\hat{\alpha} = 0$, indicates that when Residuals2 = 0 the expected value of Residuals1 is also 0. The slope $\hat{\beta} = 0.257$ shows that for every one-unit increase in Residuals2, the value of Residuals1 increases by 0.257, on average. 
	\end{enumerate}
	
	\newpage	

\section*{Question 5}
\noindent What if the incumbent's vote share is affected by both the president's popularity and the difference in spending between incumbent and challenger? 
	\begin{enumerate}
		\item Run a regression where the outcome variable is the incumbent's \texttt{voteshare} and the explanatory variables are \texttt{difflog} and \texttt{presvote}.	
		\lstinputlisting[language=R, firstline=111, lastline=113]{PS03_DP.R}
		\begin{table}[!htbp] \centering 
			\caption{} 
			\label{} 
			\begin{tabular}{@{\extracolsep{5pt}}lc} 
				\\[-1.8ex]\hline 
				\hline \\[-1.8ex] 
				& \multicolumn{1}{c}{\textit{Dependent variable:}} \\ 
				\cline{2-2} 
				\\[-1.8ex] & Incumbent Vote Share \\ 
				\hline \\[-1.8ex] 
				Campaign Spending Difference& 0.036$^{***}$ \\ 
				& (0.001) \\ 
				& \\ 
				Presidential Vote Share & 0.257$^{***}$ \\ 
				& (0.012) \\ 
				& \\ 
				Constant & 0.449$^{***}$ \\ 
				& (0.006) \\ 
				& \\ 
				\hline \\[-1.8ex] 
				\textit{Note:}  & \multicolumn{1}{r}{$^{*}$p$<$0.1; $^{**}$p$<$0.05; $^{***}$p$<$0.01} \\ 
			\end{tabular} 
		\end{table} 
		
		\item Write the prediction equation.	
		\begin{center}
		$\widehat{\text {VoteShare}} = \hat{\alpha} + \hat{\beta_1} \times \text{DiffLog} + \hat{\beta_2} \times \text{PresVote}$ \\[0.2cm]
		$\widehat{\text {VoteShare}} = 0.449 + 0.036 \times \text{DiffLog} + 0.257 \times \text{PresVote}$
		\end{center}
		The intercept $\hat{\alpha} = 0.449$, represents the expected incumbent vote share when both the difference in campaign spending between incumbent and challenger and the presidential vote share are 0. The slope $\hat{\beta_1} = 0.036$ shows that for every one-unit increase in campaign spending difference, the incumbent vote share increases by 0.036, on average, holding everything else constant. The slope $\hat{\beta_2} = 0.257$ indicates that for every one-unit increase in presidential vote share, the incumbent vote share increases by 0.257, on average, holding everything else constant. 
		\newpage
		\item What is it in this output that is identical to the output in Question 4? Why do you think this is the case? \\ [0.5cm]
		The coefficient (and standard error) of the presidential vote share (\texttt{presvote}) in this multivariate regression is identical to the coefficient (and standard error) obtained by regressing \texttt{Residuals1} on \texttt{Residuals2}. This can be explained by understanding what each of these elements represents. Specifically, \texttt{Residuals1} show how much variation in \texttt{voteshare} is unexplained by \texttt{difflog}, while  \texttt{Residuals2} indicate how much variation in \texttt{presvote} is not explained by \texttt{difflog}. \\ [0.2cm]
		Regressing \texttt{Residuals1} on \texttt{Residuals2}, we obtain the unconfounded effect of \texttt{presvote} on \texttt{voteshare}, after removing the effect of \texttt{difflog} from both variables. Within the multivariate regression, the coefficient of \texttt{presvote} also represents the effect of this variable on \texttt{voteshare}, while controlling for \texttt{difflog}, as its effect is accounted for separately. Thus, the coefficients of \texttt{presvote} and \texttt{Residuals2} are identical because they describe the same association within the data. 
	\end{enumerate}
\end{document}
